Aus $s = \frac{1}{2}g\,t^2$ folgt für die beiden Fallzeiten
\begin{align}
 t_\text{Ä} &= \taeF = \sqrt{\frac{2\cdot\sfa}{\gae}} = \tae \\
 t_\text{P} &= \tpoF = \sqrt{\frac{2\cdot\sfa}{\gpo}} = \tpo
\end{align}
Am Äquator dauert der Fall länger, um
\begin{align}
 \Delta t &= \dtF = \tae - \tpo = \dt \approx \result{\dtP{-3}{2}}
\end{align}
Relativ zur Fallzeit am Nordpol sind das
\begin{align}
 \frac{\Delta t}{t_\text{P}} &= \relF = \rel \approx \result{\relP{0}{2}}
\end{align}
Da $t \propto 1/\sqrt{g}$ ist, beträgt der relative Unterschied der Fallzeiten etwa die Hälfte des relativen Unterschieds der Fallbeschleunigungen.
